\documentclass[journal]{IEEEtran}

\usepackage{circuitikz}   
\usepackage{tikz}
\usepackage[spanish, es-noshorthands]{babel}
\usepackage{float}
\usepackage{instructivo}
\usepackage{tabularx}
\usepackage{array}
\newcolumntype{Y}{>{\raggedright\arraybackslash}X}
\graphicspath{{./}{./fig/}}
\usepackage[natbib=true,style=trad-unsrt,backend=biber]{biblatex}
\addbibresource{literatura.bib}

\hyphenation{op-tical net-works semi-conduc-tor}

\usepackage{amsmath}
\usetikzlibrary{shapes, arrows.meta, positioning, calc}
 
% --- Estilos ---
\tikzset{
  block/.style={
    draw, rectangle,
    minimum width=3.2cm, minimum height=1.2cm,
    align=center, font=\small
  },
  line/.style={draw, -{Latex[length=2.5mm, width=2mm]}},
  siglabel/.style={font=\footnotesize, text=blue!65!black},
  blklabel/.style={font=\footnotesize, right=0.15cm, text=black}
}

\begin{document}
\setlength{\parindent}{4.2mm}
\title{Controlador para convertidor buck-boost}

\author{Matías~A.~Camacho,~\IEEEmembership{Estudiante,~TEC},
        Juan~P.~Elizondo,~\IEEEmembership{Estudiante,~TEC},
        Yuliana~M.~Salas,~\IEEEmembership{Estudiante,~TEC}
        y~Daniel~E.~Salas,~\IEEEmembership{Estudiante,~TEC}
}

% The paper headers
\markboth{Laboratorio de control automático, IIS2026}%
{Laboratorio de control automático}


% make the title area
\maketitle


%%%%%%%%%%%%%%%%%%%%%%%%%%%%%%%%%%%%%%%%%%%%%%%%%%%%%%%%%%%%%%%%%%%%%%%%%%%%%%%%%%%%%%%%%%%%%%
%%%%%%%%%%%%%%%%%%%%%%%%%%%%%%%%%%%%%%%%%%%%%%%%%%%%%%%%%%%%%%%%%%%%%%%%%%%%%%%%%%%%%%%%%%%%%%
\section{Introducción}

\IEEEPARstart{E}{l} desarrollo


\printbibliography

%%%%%%%%%%%%%%%%%%%%%%%%%%%%%%%%%%%%%%%%%%%%%%%%%%%%%%%%%%%%%%%%%%%%%%%%%%%%%%%%%%%%
%%%%%%%%%%%%%%%%%%%%%%%%%%%%%%%%%%%%%%%%%%%%%%%%%%%%%%%%%%%%%%%%%%%%%%%%%%%%%%%%%%%%
\section{Biografías de los autores}

 
\begin{IEEEbiographynophoto}{Matías A. Camacho}
        Estudiante del Instituto Tecnológico de Costa Rica (TEC) en la carrera de ingeniería electrónica.
        
        Correo: jeacamacho@estudiantec.cr
\end{IEEEbiographynophoto}

\begin{IEEEbiographynophoto}{Juan P. Elizondo}
        Estudiante de Ingeniería Electrónica en el Tecnológico de Costa Rica (TEC).
        Cuenta con preparación y/o experiencia en áreas como:
        \begin{itemize}
            \item Arquitectura básica de redes, CISCO CCNA V7 (ITN), (2021).
            \item Fundamentos de ciberseguridad, CISCO Systems, (2022).
            \item Programa de tutorías estudiantiles, Tecnológico de Costa Rica, (2024).
            \item Diseño de circuitos digitales, TEC (2025).
        \end{itemize}
        
        Correo: juelizondo@estudiantec.cr
\end{IEEEbiographynophoto}

\begin{IEEEbiographynophoto}{Yuliana M. Salas}
        Estudiante del Instituto Tecnológico de Costa Rica (TEC) en la carrera de ingeniería electrónica.
        
        Correo: yulisalas0130@estudiantec.cr
\end{IEEEbiographynophoto}

\begin{IEEEbiographynophoto}{Daniel E. Salas}
        Estudiante del Instituto Tecnológico de Costa Rica (TEC) en la carrera de ingeniería electrónica.
        
        Correo: desteban2201@estudiantec.cr
\end{IEEEbiographynophoto}

\vfill

\end{document}